% Dades pròpies del llibre. Es carrega abans del preàmbul compartit.
% Una sola font de veritat per a la portada, la pàgina de crèdits i les
% metadades del PDF. Actualitzeu versió i data a cada versió publicada.
\newcommand{\llibretitol}{Una introducció a la teoria de categories}
% El subtítol només surt a la portada; a la resta de llocs n'hi ha prou amb el títol.
\newcommand{\llibresubtitol}{Formació prèvia per aprendre lògica categòrica}
\newcommand{\llibreautor}{Miquel Àngel Perelló}
\newcommand{\llibreversio}{1.7}
% Els enllaços ja inclouen \url (no s'han d'embolcallar de nou).
% Data editorial fixa (reproduïble), en lloc de \today.
\newcommand{\llibredata}{Octubre de 2026}
\newcommand{\llibreany}{2026}
\newcommand{\llibreweb}{\url{https://aprendes.com/logcat/}}
\newcommand{\llibrerepositori}{\url{https://github.com/maperello/aprendes}}
\newcommand{\llibreerrates}{\url{https://github.com/maperello/aprendes/issues}}

\title{\llibretitol\\[0.8em]\large\llibresubtitol}
\author{\llibreautor}
\date{\llibredata}

% S'ha de cridar després de carregar hyperref (és a preamble.tex).
\newcommand{\hypermetadata}{%
  \hypersetup{%
    pdftitle={\llibretitol},
    pdfauthor={\llibreautor},
    pdfsubject={Introducció a la teoria de categories orientada cap a la lògica categòrica. Versió \llibreversio\ (\llibredata). Llicència CC BY-SA 4.0},
    pdfkeywords={teoria de categories, lògica categòrica, topos, Yoneda, adjuncions}%
  }%
}
